\documentclass[border=8pt]{standalone}
\usepackage{ctex}
\usepackage{amsmath}
\usepackage{tikz}
\begin{document}
\definecolor{mmRootFill}{RGB}{18,85,126}
\definecolor{mmLOneFill}{RGB}{234,148,136}
\begin{tikzpicture}[
  mmroot/.style={draw=black!75, line width=1.0pt, rounded corners=3.5pt, fill=mmRootFill, inner sep=5pt, font=\normalsize\bfseries},
  mml1/.style={draw=black!55, line width=0.9pt, rounded corners=2.5pt, fill=mmLOneFill, inner sep=4.5pt, font=\small\bfseries, align=left, text width=4.8cm},
  mml2/.style={draw=black!45, line width=0.8pt, rounded corners=2pt, fill=white, inner sep=4pt, font=\small, align=left, text width=6.2cm},
]
  %% 连线：根 → 一级 → 二级（节点填充不透明，盖住线头）
  \draw[black!55, line width=0.9pt] (4.35, 3.45) -- (4.8, 7.475);
  \draw[black!55, line width=0.9pt] (4.8, 7.475) -- (10.5, 7.475);
  \draw[black!55, line width=0.9pt] (4.35, 3.45) -- (4.8, 5.175);
  \draw[black!55, line width=0.9pt] (4.8, 5.175) -- (10.5, 6.325);
  \draw[black!55, line width=0.9pt] (4.8, 5.175) -- (10.5, 5.175);
  \draw[black!55, line width=0.9pt] (4.8, 5.175) -- (10.5, 4.025);
  \draw[black!55, line width=0.9pt] (4.35, 3.45) -- (4.8, 1.725);
  \draw[black!55, line width=0.9pt] (4.8, 1.725) -- (10.5, 2.875);
  \draw[black!55, line width=0.9pt] (4.8, 1.725) -- (10.5, 1.725);
  \draw[black!55, line width=0.9pt] (4.8, 1.725) -- (10.5, 0.575);
  \draw[black!55, line width=0.9pt] (4.35, 3.45) -- (4.8, -0.575);
  \draw[black!55, line width=0.9pt] (4.8, -0.575) -- (10.5, -0.575);
  %% 节点
  \node[mmroot, anchor=east] at (4.35, 3.45) {2.3 圆及其方程};
  \node[mml1, anchor=west] at (4.6, 7.475) {2.3.1 圆的标准方程};
  \node[mml2, anchor=west] at (10.3, 7.475) {圆的标准方程};
  \node[mml1, anchor=west] at (4.6, 5.175) {2.3.2 圆的一般方程};
  \node[mml2, anchor=west] at (10.3, 6.325) {圆的一般方程的概念};
  \node[mml2, anchor=west] at (10.3, 5.175) {圆的一般方程对应的圆心和半径};
  \node[mml2, anchor=west] at (10.3, 4.025) {圆系方程};
  \node[mml1, anchor=west] at (4.6, 1.725) {2.3.3 直线与圆的位置关系};
  \node[mml2, anchor=west] at (10.3, 2.875) {直线与圆的位置关系及判断};
  \node[mml2, anchor=west] at (10.3, 1.725) {阿波罗尼斯圆};
  \node[mml2, anchor=west] at (10.3, 0.575) {隐圆的定义};
  \node[mml1, anchor=west] at (4.6, -0.575) {2.3.4 圆与圆的位置关系};
  \node[mml2, anchor=west] at (10.3, -0.575) {圆与圆位置关系的判定};
\end{tikzpicture}
\end{document}
